\documentclass[11pt,a4paper]{article}
% lesson-preamble.tex -- Shared preamble for all Phase 00 lessons
% Usage: % lesson-preamble.tex -- Shared preamble for all Phase 00 lessons
% Usage: % lesson-preamble.tex -- Shared preamble for all Phase 00 lessons
% Usage: \input{../lesson-preamble} (from subdomain directories)

% ---- Encoding and fonts ----
\usepackage[utf8]{inputenc}
\usepackage[T1]{fontenc}

% ---- Mathematics ----
\usepackage{amsmath,amssymb,amsthm}
\usepackage{mathtools}

% ---- Page geometry ----
\usepackage[margin=1in,headheight=14pt]{geometry}

% ---- Hyperlinks ----
\usepackage[colorlinks=true,linkcolor=red,citecolor=blue,urlcolor=blue]{hyperref}

% ---- Lists ----
\usepackage{enumitem}

% ---- Colors and boxes ----
\usepackage{xcolor}
\usepackage[theorems,skins,breakable]{tcolorbox}

% ---- Header/Footer ----
\usepackage{fancyhdr}
\renewcommand{\headrulewidth}{0.4pt}

% ---- Tables ----
\usepackage{booktabs}

% ---- Bibliography ----
\usepackage{natbib}
\bibliographystyle{plainnat}

% --------------------------------------------------------------------------
% Theorem environments
% --------------------------------------------------------------------------
\theoremstyle{plain}
\newtheorem{theorem}{Theorem}[section]
\newtheorem{lemma}[theorem]{Lemma}
\newtheorem{proposition}[theorem]{Proposition}
\newtheorem{corollary}[theorem]{Corollary}

\theoremstyle{definition}
\newtheorem{definition}[theorem]{Definition}
\newtheorem{example}[theorem]{Example}
\newtheorem{exercise}[theorem]{Exercise}
\newtheorem{assumption}[theorem]{Assumption}

\theoremstyle{remark}
\newtheorem{remark}[theorem]{Remark}

% --------------------------------------------------------------------------
% Colored boxes
% --------------------------------------------------------------------------
\newtcolorbox{keyresult}[1][]{%
  colback=green!5!white,
  colframe=green!50!black,
  fonttitle=\bfseries,
  title={Key Result},
  breakable,
  #1
}

\newtcolorbox{rlconnection}[1][]{%
  colback=blue!5!white,
  colframe=blue!50!black,
  fonttitle=\bfseries,
  title={RL/ML Connection},
  breakable,
  #1
}

\newtcolorbox{warning}[1][]{%
  colback=red!5!white,
  colframe=red!50!black,
  fonttitle=\bfseries,
  title={Warning},
  breakable,
  #1
}

% --------------------------------------------------------------------------
% Custom commands -- number sets
% --------------------------------------------------------------------------
\newcommand{\R}{\mathbb{R}}
\newcommand{\N}{\mathbb{N}}
\newcommand{\Q}{\mathbb{Q}}
\newcommand{\Z}{\mathbb{Z}}
\newcommand{\C}{\mathbb{C}}
\newcommand{\F}{\mathbb{F}}

% --------------------------------------------------------------------------
% Custom commands -- probability and expectation
% --------------------------------------------------------------------------
\newcommand{\E}{\mathbb{E}}
\newcommand{\Prob}{\mathbb{P}}
\newcommand{\Var}{\operatorname{Var}}
\newcommand{\Cov}{\operatorname{Cov}}
\newcommand{\indicator}{\mathbf{1}}

% --------------------------------------------------------------------------
% Custom commands -- convergence arrows
% --------------------------------------------------------------------------
\newcommand{\cas}{\xrightarrow{\text{a.s.}}}
\newcommand{\cprob}{\xrightarrow{\Prob}}
\newcommand{\clp}[1]{\xrightarrow{L^{#1}}}
\newcommand{\cdist}{\xrightarrow{d}}

% --------------------------------------------------------------------------
% Custom commands -- common operators
% --------------------------------------------------------------------------
\DeclareMathOperator{\dom}{dom}
\DeclareMathOperator{\epi}{epi}
\DeclareMathOperator{\conv}{conv}
\DeclareMathOperator{\relint}{relint}
\DeclareMathOperator{\aff}{aff}
\DeclareMathOperator{\cl}{cl}
\DeclareMathOperator{\interior}{int}
\DeclareMathOperator{\tr}{tr}
\DeclareMathOperator{\diag}{diag}
\DeclareMathOperator{\rank}{rank}
\DeclareMathOperator{\sgn}{sgn}
\DeclareMathOperator{\supp}{supp}
\DeclareMathOperator{\argmin}{arg\,min}
\DeclareMathOperator{\argmax}{arg\,max}
\DeclareMathOperator{\prox}{prox}
\DeclareMathOperator{\dist}{dist}
\DeclareMathOperator{\diam}{diam}
\DeclareMathOperator{\Span}{span}

% --------------------------------------------------------------------------
% Custom commands -- linear algebra operators
% --------------------------------------------------------------------------
\DeclareMathOperator{\nullity}{nullity}
\DeclareMathOperator{\range}{range}
\DeclareMathOperator{\nullspace}{null}
\DeclareMathOperator{\proj}{proj}

% --------------------------------------------------------------------------
% Custom commands -- measure theory
% --------------------------------------------------------------------------
\newcommand{\powerset}{\mathcal{P}}
\newcommand{\Borel}{\mathcal{B}}
 (from subdomain directories)

% ---- Encoding and fonts ----
\usepackage[utf8]{inputenc}
\usepackage[T1]{fontenc}

% ---- Mathematics ----
\usepackage{amsmath,amssymb,amsthm}
\usepackage{mathtools}

% ---- Page geometry ----
\usepackage[margin=1in,headheight=14pt]{geometry}

% ---- Hyperlinks ----
\usepackage[colorlinks=true,linkcolor=red,citecolor=blue,urlcolor=blue]{hyperref}

% ---- Lists ----
\usepackage{enumitem}

% ---- Colors and boxes ----
\usepackage{xcolor}
\usepackage[theorems,skins,breakable]{tcolorbox}

% ---- Header/Footer ----
\usepackage{fancyhdr}
\renewcommand{\headrulewidth}{0.4pt}

% ---- Tables ----
\usepackage{booktabs}

% ---- Bibliography ----
\usepackage{natbib}
\bibliographystyle{plainnat}

% --------------------------------------------------------------------------
% Theorem environments
% --------------------------------------------------------------------------
\theoremstyle{plain}
\newtheorem{theorem}{Theorem}[section]
\newtheorem{lemma}[theorem]{Lemma}
\newtheorem{proposition}[theorem]{Proposition}
\newtheorem{corollary}[theorem]{Corollary}

\theoremstyle{definition}
\newtheorem{definition}[theorem]{Definition}
\newtheorem{example}[theorem]{Example}
\newtheorem{exercise}[theorem]{Exercise}
\newtheorem{assumption}[theorem]{Assumption}

\theoremstyle{remark}
\newtheorem{remark}[theorem]{Remark}

% --------------------------------------------------------------------------
% Colored boxes
% --------------------------------------------------------------------------
\newtcolorbox{keyresult}[1][]{%
  colback=green!5!white,
  colframe=green!50!black,
  fonttitle=\bfseries,
  title={Key Result},
  breakable,
  #1
}

\newtcolorbox{rlconnection}[1][]{%
  colback=blue!5!white,
  colframe=blue!50!black,
  fonttitle=\bfseries,
  title={RL/ML Connection},
  breakable,
  #1
}

\newtcolorbox{warning}[1][]{%
  colback=red!5!white,
  colframe=red!50!black,
  fonttitle=\bfseries,
  title={Warning},
  breakable,
  #1
}

% --------------------------------------------------------------------------
% Custom commands -- number sets
% --------------------------------------------------------------------------
\newcommand{\R}{\mathbb{R}}
\newcommand{\N}{\mathbb{N}}
\newcommand{\Q}{\mathbb{Q}}
\newcommand{\Z}{\mathbb{Z}}
\newcommand{\C}{\mathbb{C}}
\newcommand{\F}{\mathbb{F}}

% --------------------------------------------------------------------------
% Custom commands -- probability and expectation
% --------------------------------------------------------------------------
\newcommand{\E}{\mathbb{E}}
\newcommand{\Prob}{\mathbb{P}}
\newcommand{\Var}{\operatorname{Var}}
\newcommand{\Cov}{\operatorname{Cov}}
\newcommand{\indicator}{\mathbf{1}}

% --------------------------------------------------------------------------
% Custom commands -- convergence arrows
% --------------------------------------------------------------------------
\newcommand{\cas}{\xrightarrow{\text{a.s.}}}
\newcommand{\cprob}{\xrightarrow{\Prob}}
\newcommand{\clp}[1]{\xrightarrow{L^{#1}}}
\newcommand{\cdist}{\xrightarrow{d}}

% --------------------------------------------------------------------------
% Custom commands -- common operators
% --------------------------------------------------------------------------
\DeclareMathOperator{\dom}{dom}
\DeclareMathOperator{\epi}{epi}
\DeclareMathOperator{\conv}{conv}
\DeclareMathOperator{\relint}{relint}
\DeclareMathOperator{\aff}{aff}
\DeclareMathOperator{\cl}{cl}
\DeclareMathOperator{\interior}{int}
\DeclareMathOperator{\tr}{tr}
\DeclareMathOperator{\diag}{diag}
\DeclareMathOperator{\rank}{rank}
\DeclareMathOperator{\sgn}{sgn}
\DeclareMathOperator{\supp}{supp}
\DeclareMathOperator{\argmin}{arg\,min}
\DeclareMathOperator{\argmax}{arg\,max}
\DeclareMathOperator{\prox}{prox}
\DeclareMathOperator{\dist}{dist}
\DeclareMathOperator{\diam}{diam}
\DeclareMathOperator{\Span}{span}

% --------------------------------------------------------------------------
% Custom commands -- linear algebra operators
% --------------------------------------------------------------------------
\DeclareMathOperator{\nullity}{nullity}
\DeclareMathOperator{\range}{range}
\DeclareMathOperator{\nullspace}{null}
\DeclareMathOperator{\proj}{proj}

% --------------------------------------------------------------------------
% Custom commands -- measure theory
% --------------------------------------------------------------------------
\newcommand{\powerset}{\mathcal{P}}
\newcommand{\Borel}{\mathcal{B}}
 (from subdomain directories)

% ---- Encoding and fonts ----
\usepackage[utf8]{inputenc}
\usepackage[T1]{fontenc}

% ---- Mathematics ----
\usepackage{amsmath,amssymb,amsthm}
\usepackage{mathtools}

% ---- Page geometry ----
\usepackage[margin=1in,headheight=14pt]{geometry}

% ---- Hyperlinks ----
\usepackage[colorlinks=true,linkcolor=red,citecolor=blue,urlcolor=blue]{hyperref}

% ---- Lists ----
\usepackage{enumitem}

% ---- Colors and boxes ----
\usepackage{xcolor}
\usepackage[theorems,skins,breakable]{tcolorbox}

% ---- Header/Footer ----
\usepackage{fancyhdr}
\renewcommand{\headrulewidth}{0.4pt}

% ---- Tables ----
\usepackage{booktabs}

% ---- Bibliography ----
\usepackage{natbib}
\bibliographystyle{plainnat}

% --------------------------------------------------------------------------
% Theorem environments
% --------------------------------------------------------------------------
\theoremstyle{plain}
\newtheorem{theorem}{Theorem}[section]
\newtheorem{lemma}[theorem]{Lemma}
\newtheorem{proposition}[theorem]{Proposition}
\newtheorem{corollary}[theorem]{Corollary}

\theoremstyle{definition}
\newtheorem{definition}[theorem]{Definition}
\newtheorem{example}[theorem]{Example}
\newtheorem{exercise}[theorem]{Exercise}
\newtheorem{assumption}[theorem]{Assumption}

\theoremstyle{remark}
\newtheorem{remark}[theorem]{Remark}

% --------------------------------------------------------------------------
% Colored boxes
% --------------------------------------------------------------------------
\newtcolorbox{keyresult}[1][]{%
  colback=green!5!white,
  colframe=green!50!black,
  fonttitle=\bfseries,
  title={Key Result},
  breakable,
  #1
}

\newtcolorbox{rlconnection}[1][]{%
  colback=blue!5!white,
  colframe=blue!50!black,
  fonttitle=\bfseries,
  title={RL/ML Connection},
  breakable,
  #1
}

\newtcolorbox{warning}[1][]{%
  colback=red!5!white,
  colframe=red!50!black,
  fonttitle=\bfseries,
  title={Warning},
  breakable,
  #1
}

% --------------------------------------------------------------------------
% Custom commands -- number sets
% --------------------------------------------------------------------------
\newcommand{\R}{\mathbb{R}}
\newcommand{\N}{\mathbb{N}}
\newcommand{\Q}{\mathbb{Q}}
\newcommand{\Z}{\mathbb{Z}}
\newcommand{\C}{\mathbb{C}}
\newcommand{\F}{\mathbb{F}}

% --------------------------------------------------------------------------
% Custom commands -- probability and expectation
% --------------------------------------------------------------------------
\newcommand{\E}{\mathbb{E}}
\newcommand{\Prob}{\mathbb{P}}
\newcommand{\Var}{\operatorname{Var}}
\newcommand{\Cov}{\operatorname{Cov}}
\newcommand{\indicator}{\mathbf{1}}

% --------------------------------------------------------------------------
% Custom commands -- convergence arrows
% --------------------------------------------------------------------------
\newcommand{\cas}{\xrightarrow{\text{a.s.}}}
\newcommand{\cprob}{\xrightarrow{\Prob}}
\newcommand{\clp}[1]{\xrightarrow{L^{#1}}}
\newcommand{\cdist}{\xrightarrow{d}}

% --------------------------------------------------------------------------
% Custom commands -- common operators
% --------------------------------------------------------------------------
\DeclareMathOperator{\dom}{dom}
\DeclareMathOperator{\epi}{epi}
\DeclareMathOperator{\conv}{conv}
\DeclareMathOperator{\relint}{relint}
\DeclareMathOperator{\aff}{aff}
\DeclareMathOperator{\cl}{cl}
\DeclareMathOperator{\interior}{int}
\DeclareMathOperator{\tr}{tr}
\DeclareMathOperator{\diag}{diag}
\DeclareMathOperator{\rank}{rank}
\DeclareMathOperator{\sgn}{sgn}
\DeclareMathOperator{\supp}{supp}
\DeclareMathOperator{\argmin}{arg\,min}
\DeclareMathOperator{\argmax}{arg\,max}
\DeclareMathOperator{\prox}{prox}
\DeclareMathOperator{\dist}{dist}
\DeclareMathOperator{\diam}{diam}
\DeclareMathOperator{\Span}{span}

% --------------------------------------------------------------------------
% Custom commands -- linear algebra operators
% --------------------------------------------------------------------------
\DeclareMathOperator{\nullity}{nullity}
\DeclareMathOperator{\range}{range}
\DeclareMathOperator{\nullspace}{null}
\DeclareMathOperator{\proj}{proj}

% --------------------------------------------------------------------------
% Custom commands -- measure theory
% --------------------------------------------------------------------------
\newcommand{\powerset}{\mathcal{P}}
\newcommand{\Borel}{\mathcal{B}}

\pagestyle{fancy}
\fancyhf{}
\fancyhead[L]{\textsc{Lesson 11: Expectation, Variance, and Joint Distributions}}
\fancyhead[R]{\thepage}
\title{Lesson 11: Expectation, Variance, and Joint Distributions}
\author{AI/ML Mathematics}
\date{}
\begin{document}
\maketitle
\tableofcontents
\newpage

%% ============================================================
%% SECTION 1: MOTIVATION AND OVERVIEW
%% ============================================================

\section{Motivation and Overview}
\label{sec:motivation}

An RL agent estimates the value of a state as an expected sum of future
rewards: $V^\pi(s) = \E\bigl[\sum_{t=0}^\infty \gamma^t R_t \mid
S_0 = s\bigr]$. This expression involves an expectation over many random
variables (the rewards $R_t$), decomposes as a sum of individual
expectations by linearity, and connects values at successive time steps
through the tower property
$\E[V^\pi(S_{t+1}) \mid S_t] = \E[\E[\sum \gamma^k R_{t+1+k} \mid
S_{t+1}] \mid S_t]$. To make such manipulations rigorous, we need
a systematic treatment of expectation, variance, covariance, and the way
multiple random variables interact through joint distributions. How does
the expectation of a sum relate to the sum of expectations? When does the
variance of a sum simplify? How do we condition on one random variable to
compute the distribution of another?

The tools developed in this lesson will let us prove the linearity of
expectation without assuming independence, establish the variance
decomposition for sums, define joint and conditional distributions, and
derive the tower property of conditional expectation -- the algebraic
engine behind the Bellman equation.

\begin{rlconnection}[title={Why Expectation and Joint Distributions
    Matter for RL/ML}]
\begin{itemize}[nosep]
    \item \textbf{Value function decomposition}: The Bellman equation
    $V^\pi(s) = \E[R + \gamma V^\pi(S') \mid S = s]$ uses linearity of
    expectation to split a value function into an immediate reward and a
    discounted future value.
    \item \textbf{Bellman equation via the tower property}: The recursive
    structure of value functions follows from the law of total
    expectation (tower property), which allows conditioning on the next
    state.
    \item \textbf{REINFORCE variance reduction}: The baseline subtraction
    $\nabla J \approx \E[\nabla \log\pi \cdot (G - b)]$ uses the fact
    that adding a baseline $b(s)$ does not change the expectation (by
    linearity) but can reduce variance (by the variance of a
    difference formula).
\end{itemize}
\end{rlconnection}

\paragraph{Prerequisites.}
Random variables, probability mass functions, probability density
functions, cumulative distribution functions, and common distributions
(Lesson~10). Conditional probability and independence of events
(Lesson~9).

\paragraph{Learning objectives.}
After completing this lesson, the student will be able to:
\begin{itemize}
    \item \emph{State} and \emph{prove} the linearity of expectation for
    both discrete and continuous random variables, without assuming
    independence.
    \item \emph{Compute} expectations using the Law of the Unconscious
    Statistician (LOTUS).
    \item \emph{Derive} the computational formula
    $\Var(X) = \E[X^2] - (\E[X])^2$ and the affine transformation rule
    $\Var(aX + b) = a^2\Var(X)$.
    \item \emph{Define} covariance and correlation, and \emph{prove}
    the variance of a sum formula.
    \item \emph{Construct} an example showing that uncorrelated random
    variables need not be independent.
    \item \emph{Define} joint, marginal, and conditional distributions
    and \emph{prove} the law of total expectation (tower property) and
    the law of total variance.
\end{itemize}

%% ============================================================
%% SECTION 2: REFERENCES
%% ============================================================

\section{References}
\label{sec:references}

\begin{itemize}
    \item \citet[Chapters 4, 6--7]{ross2019} -- Expectation, variance,
    jointly distributed random variables, conditional distributions.
    \item \citet[Chapters 4, 6--7]{blitzstein2019} -- Expectation,
    moments, joint distributions, conditional expectation with many
    examples from applied probability.
    \item \citet[Chapter 3]{grimmett2020} -- Discrete and continuous
    random variables, expectation, variance, generating functions.
    \item \citet[Chapters 2--4]{bertsekas2008} -- Expectation, variance,
    multiple random variables, conditioning.
    \item \citet[Sections 1.5--1.6]{durrett2019} -- Expectation and
    conditional expectation in the measure-theoretic framework.
\end{itemize}

%% ============================================================
%% SECTION 3: EXPECTATION
%% ============================================================

\section{Expectation}
\label{sec:expectation}

We develop the properties of expectation for discrete random variables
first, then extend to the continuous case.

\subsection{Expectation for Discrete Random Variables}

Recall from Lesson~10 (Definition~4.2) that if $X$ is a discrete random
variable with PMF $p_X$, then $\E[X] = \sum_x x\,p_X(x)$ whenever the
sum converges absolutely.

\begin{theorem}[Law of the Unconscious Statistician -- Discrete Case]
\label{thm:lotus_discrete}
Let $X$ be a discrete random variable and $g \colon \R \to \R$ any
function. Then
\[
    \E[g(X)] = \sum_x g(x) \, p_X(x),
\]
provided the sum converges absolutely.
\end{theorem}

\begin{proof}
Define $Y = g(X)$. The random variable $Y$ takes values in
$\{g(x) : x \in \mathrm{range}(X)\}$. For each value $y$ in the range
of $Y$,
\[
    p_Y(y) = \Prob(Y = y) = \Prob(g(X) = y)
    = \sum_{\{x : g(x) = y\}} p_X(x).
\]
Therefore
\[
    \E[Y] = \sum_y y \, p_Y(y)
    = \sum_y y \sum_{\{x : g(x) = y\}} p_X(x)
    = \sum_y \sum_{\{x : g(x) = y\}} g(x) \, p_X(x)
    = \sum_x g(x) \, p_X(x),
\]
where the last equality rearranges the double sum: each $x$ in the range
of $X$ appears exactly once, under the unique $y = g(x)$.
\end{proof}

\begin{keyresult}[title={Linearity of Expectation}]
Linearity of expectation is arguably the single most useful property
in probability. It states that $\E[aX + bY] = a\E[X] + b\E[Y]$ for
\emph{any} random variables $X$ and $Y$ (with finite expectations) and
constants $a, b \in \R$. Crucially, no independence assumption is
needed. This property underlies the decomposition of value functions,
the analysis of Monte Carlo estimators, and virtually every expected
reward computation in RL.
\end{keyresult}

\begin{theorem}[Linearity of Expectation -- Discrete Case]
\label{thm:linearity_discrete}
Let $X$ and $Y$ be discrete random variables defined on the same
probability space with finite expectations. Then for any
$a, b \in \R$,
\[
    \E[aX + bY] = a\E[X] + b\E[Y].
\]
\end{theorem}

\begin{proof}
Let $p_{X,Y}(x,y) = \Prob(X = x, Y = y)$ be the joint PMF. By the
definition of expectation applied to the random variable $aX + bY$,
\begin{align*}
    \E[aX + bY]
    &= \sum_x \sum_y (ax + by) \, p_{X,Y}(x,y) \\
    &= a \sum_x \sum_y x \, p_{X,Y}(x,y)
       + b \sum_x \sum_y y \, p_{X,Y}(x,y).
\end{align*}
For the first term, note that
$\sum_y p_{X,Y}(x,y) = p_X(x)$ (the marginal PMF of $X$), so
\[
    \sum_x \sum_y x \, p_{X,Y}(x,y)
    = \sum_x x \sum_y p_{X,Y}(x,y)
    = \sum_x x \, p_X(x)
    = \E[X].
\]
Similarly, $\sum_x p_{X,Y}(x,y) = p_Y(y)$, so the second term
equals $\E[Y]$. Combining:
$\E[aX + bY] = a\E[X] + b\E[Y]$.
\end{proof}

\begin{remark}
\label{rem:linearity_general}
The proof extends immediately to any finite sum: for random variables
$X_1, \ldots, X_n$ and constants $a_1, \ldots, a_n$,
\[
    \E\!\left[\sum_{i=1}^n a_i X_i\right]
    = \sum_{i=1}^n a_i \E[X_i].
\]
This is proved by induction on $n$, applying
Theorem~\ref{thm:linearity_discrete} at each step.
\end{remark}

\begin{example}[Indicator Expectation]
\label{ex:indicator_expectation}
Let $A$ be an event and $\indicator_A$ the indicator random variable
with $\indicator_A(\omega) = 1$ if $\omega \in A$ and $0$ otherwise.
Then $\indicator_A \sim \mathrm{Bernoulli}(\Prob(A))$, so
$\E[\indicator_A] = \Prob(A)$. This links probability to expectation:
every probability is the expectation of an indicator.
\end{example}

\begin{example}[Expected Number of Fixed Points]
\label{ex:fixed_points}
A random permutation $\sigma$ of $\{1, 2, \ldots, n\}$ is chosen
uniformly. Let $X$ be the number of fixed points ($\sigma(i) = i$).
Write $X = \sum_{i=1}^n \indicator_{\{\sigma(i)=i\}}$. Each indicator
has expectation $\Prob(\sigma(i) = i) = 1/n$ (since each element is
equally likely to land in any position). By linearity,
$\E[X] = n \cdot (1/n) = 1$,
regardless of $n$. Note that the indicators are \emph{not} independent,
but linearity of expectation applies nonetheless.
\end{example}

\subsection{Expectation for Continuous Random Variables}

\begin{theorem}[LOTUS and Linearity -- Continuous Case]
\label{thm:lotus_continuous}
Let $X$ be a continuous random variable with PDF $f_X$ and
$g \colon \R \to \R$. Then
$\E[g(X)] = \int_{-\infty}^{\infty} g(x) f_X(x)\,dx$,
provided the integral converges absolutely. In particular, linearity
holds: for jointly continuous $X, Y$ with joint PDF $f_{X,Y}$,
\[
    \E[aX + bY] = a\E[X] + b\E[Y].
\]
\end{theorem}

\begin{proof}
The LOTUS formula is analogous to the discrete case
(Theorem~\ref{thm:lotus_discrete}): one integrates $g(x) f_X(x)$
instead of summing $g(x) p_X(x)$.

For linearity,
\begin{align*}
    \E[aX + bY]
    &= \int_{-\infty}^{\infty}\int_{-\infty}^{\infty}
    (ax + by) f_{X,Y}(x,y)\,dx\,dy \\
    &= a\int_{-\infty}^{\infty}\int_{-\infty}^{\infty}
    x\,f_{X,Y}(x,y)\,dx\,dy
    + b\int_{-\infty}^{\infty}\int_{-\infty}^{\infty}
    y\,f_{X,Y}(x,y)\,dx\,dy \\
    &= a\int_{-\infty}^{\infty} x\,f_X(x)\,dx
    + b\int_{-\infty}^{\infty} y\,f_Y(y)\,dy
    = a\E[X] + b\E[Y],
\end{align*}
where $f_X(x) = \int f_{X,Y}(x,y)\,dy$ is the marginal PDF of $X$.
\end{proof}

\begin{rlconnection}[title={Linearity and Value Function Decomposition}]
The Bellman equation for a policy $\pi$ in a discounted MDP is
\[
    V^\pi(s) = \E\bigl[R_t + \gamma V^\pi(S_{t+1}) \mid S_t = s\bigr]
    = \E[R_t \mid S_t = s] + \gamma \, \E[V^\pi(S_{t+1}) \mid S_t = s].
\]
The second equality is a direct application of linearity: the
expectation of the sum (immediate reward plus discounted future value)
equals the sum of expectations. No independence between $R_t$ and
$S_{t+1}$ is assumed -- they may be highly correlated (e.g., higher
rewards might occur in states with better future prospects).
\end{rlconnection}

We now turn to variance and its properties.

%% ============================================================
%% SECTION 4: VARIANCE AND COVARIANCE
%% ============================================================

\section{Variance and Covariance}
\label{sec:variance}

\subsection{Variance Properties}

Recall from Lesson~10 that $\Var(X) = \E[(X - \E[X])^2]$ and the
computational formula $\Var(X) = \E[X^2] - (\E[X])^2$
(Lesson~10, Proposition~4.4).

\begin{proposition}[Affine Transformation of Variance]
\label{prop:affine_var}
For any random variable $X$ with finite variance and constants
$a, b \in \R$,
\[
    \Var(aX + b) = a^2 \Var(X).
\]
In particular, adding a constant does not change the variance, and
scaling by $a$ scales the variance by $a^2$.
\end{proposition}

\begin{proof}
Since $\E[aX + b] = a\E[X] + b$ by linearity,
\begin{align*}
    \Var(aX + b) &= \E\bigl[(aX + b - a\E[X] - b)^2\bigr]
    = \E\bigl[a^2(X - \E[X])^2\bigr]
    = a^2 \E\bigl[(X - \E[X])^2\bigr]
    = a^2 \Var(X). \qedhere
\end{align*}
\end{proof}

\begin{example}[Standardization]
\label{ex:standardization}
If $\E[X] = \mu$ and $\Var(X) = \sigma^2 > 0$, define the
\emph{standardized} variable $Z = (X - \mu)/\sigma$. Then
$\E[Z] = 0$ and $\Var(Z) = \Var(X)/\sigma^2 = 1$.
\end{example}

\subsection{Covariance}

\begin{definition}[Covariance]
\label{def:covariance}
The \emph{covariance} of random variables $X$ and $Y$ (with finite
second moments) is
\[
    \Cov(X, Y) = \E\bigl[(X - \E[X])(Y - \E[Y])\bigr].
\]
\end{definition}

\begin{proposition}[Computational Formula for Covariance]
\label{prop:cov_formula}
\[
    \Cov(X, Y) = \E[XY] - \E[X]\E[Y].
\]
\end{proposition}

\begin{proof}
Let $\mu_X = \E[X]$ and $\mu_Y = \E[Y]$. Then
\begin{align*}
    \Cov(X,Y) &= \E[(X - \mu_X)(Y - \mu_Y)]
    = \E[XY - \mu_X Y - X \mu_Y + \mu_X \mu_Y] \\
    &= \E[XY] - \mu_X \E[Y] - \E[X] \mu_Y + \mu_X \mu_Y
    = \E[XY] - \mu_X \mu_Y. \qedhere
\end{align*}
\end{proof}

\begin{proposition}[Properties of Covariance]
\label{prop:cov_properties}
For random variables $X, Y, Z$ with finite second moments and constants
$a, b \in \R$:
\begin{enumerate}[nosep,label=(\roman*)]
    \item \textbf{Symmetry:} $\Cov(X, Y) = \Cov(Y, X)$.
    \item \textbf{Self-covariance:} $\Cov(X, X) = \Var(X)$.
    \item \textbf{Bilinearity:}
    $\Cov(aX + bY, Z) = a\,\Cov(X, Z) + b\,\Cov(Y, Z)$.
    \item \textbf{Constant:} $\Cov(X, c) = 0$ for any constant $c$.
\end{enumerate}
\end{proposition}

\begin{proof}
\textbf{(i)} By commutativity of multiplication:
$(X - \E[X])(Y - \E[Y]) = (Y - \E[Y])(X - \E[X])$.

\textbf{(ii)} $\Cov(X, X) = \E[(X - \E[X])^2] = \Var(X)$.

\textbf{(iii)} Using the computational formula:
\begin{align*}
    \Cov(aX + bY, Z)
    &= \E[(aX + bY)Z] - \E[aX + bY]\E[Z] \\
    &= a\E[XZ] + b\E[YZ] - (a\E[X] + b\E[Y])\E[Z] \\
    &= a(\E[XZ] - \E[X]\E[Z]) + b(\E[YZ] - \E[Y]\E[Z]) \\
    &= a\,\Cov(X,Z) + b\,\Cov(Y,Z).
\end{align*}

\textbf{(iv)} $\Cov(X, c) = \E[Xc] - \E[X]\E[c]
= c\E[X] - c\E[X] = 0$.
\end{proof}

\begin{theorem}[Variance of a Sum]
\label{thm:var_sum}
For random variables $X_1, \ldots, X_n$ with finite second moments,
\[
    \Var\!\left(\sum_{i=1}^n X_i\right)
    = \sum_{i=1}^n \Var(X_i)
    + 2 \sum_{1 \leq i < j \leq n} \Cov(X_i, X_j).
\]
\end{theorem}

\begin{proof}
Let $S = \sum_{i=1}^n X_i$ and $\mu_i = \E[X_i]$. Then
$\E[S] = \sum_i \mu_i$ and
\begin{align*}
    \Var(S) &= \E\!\left[\left(\sum_{i=1}^n (X_i - \mu_i)\right)^2\right]
    = \E\!\left[\sum_{i=1}^n \sum_{j=1}^n
    (X_i - \mu_i)(X_j - \mu_j)\right] \\
    &= \sum_{i=1}^n \sum_{j=1}^n
    \E\bigl[(X_i - \mu_i)(X_j - \mu_j)\bigr]
    = \sum_{i=1}^n \sum_{j=1}^n \Cov(X_i, X_j).
\end{align*}
Separating diagonal terms ($i = j$) from off-diagonal terms ($i \neq j$)
and using symmetry $\Cov(X_i, X_j) = \Cov(X_j, X_i)$:
\[
    \Var(S) = \sum_{i=1}^n \Var(X_i)
    + 2\sum_{1 \leq i < j \leq n} \Cov(X_i, X_j). \qedhere
\]
\end{proof}

\begin{corollary}[Variance of Independent Sum]
\label{cor:var_indep_sum}
If $X_1, \ldots, X_n$ are pairwise independent (i.e.,
$\Cov(X_i, X_j) = 0$ for $i \neq j$; see
Proposition~\ref{prop:indep_uncorrelated} below), then
\[
    \Var\!\left(\sum_{i=1}^n X_i\right)
    = \sum_{i=1}^n \Var(X_i).
\]
\end{corollary}

\begin{proposition}[Independent Implies Uncorrelated]
\label{prop:indep_uncorrelated}
If $X$ and $Y$ are independent random variables with finite second
moments, then $\Cov(X, Y) = 0$.
\end{proposition}

\begin{proof}
Independence of $X$ and $Y$ implies $\E[XY] = \E[X]\E[Y]$. (For
discrete random variables: $\E[XY] = \sum_x \sum_y xy \, p_{X,Y}(x,y)
= \sum_x \sum_y xy \, p_X(x) p_Y(y) = (\sum_x x \, p_X(x))(\sum_y y
\, p_Y(y)) = \E[X]\E[Y]$; the continuous case is analogous with
integrals and the factorization $f_{X,Y}(x,y) = f_X(x) f_Y(y)$.)
By the computational formula,
$\Cov(X,Y) = \E[XY] - \E[X]\E[Y] = 0$.
\end{proof}

\begin{warning}[title={Uncorrelated Does Not Imply Independent}]
The converse of Proposition~\ref{prop:indep_uncorrelated} is false:
$\Cov(X, Y) = 0$ does not imply that $X$ and $Y$ are independent.
Covariance measures only \emph{linear} association. Two variables can be
functionally related yet uncorrelated.
\end{warning}

\begin{example}[Uncorrelated but Dependent]
\label{ex:uncorrelated_dependent}
Let $X \sim \mathrm{Uniform}(-1, 1)$ and $Y = X^2$. Then $Y$ is
completely determined by $X$ (maximal dependence). However,
$\E[X] = 0$, so $\Cov(X, Y) = \E[XY] - \E[X]\E[Y] = \E[X^3]$.
Since $x^3 f_X(x) = x^3/(2)$ is an odd function on $(-1, 1)$,
$\E[X^3] = \int_{-1}^{1} x^3/(2)\,dx = 0$.
Thus $\Cov(X, Y) = 0$ even though $Y$ is a deterministic function of $X$.
\end{example}

\subsection{Correlation Coefficient}

\begin{definition}[Correlation Coefficient]
\label{def:correlation}
The \emph{(Pearson) correlation coefficient} of $X$ and $Y$ (with
$\Var(X) > 0$ and $\Var(Y) > 0$) is
\[
    \rho(X, Y) = \frac{\Cov(X, Y)}{\sqrt{\Var(X) \Var(Y)}}.
\]
\end{definition}

\begin{proposition}
\label{prop:correlation_bounds}
$-1 \leq \rho(X, Y) \leq 1$, with $|\rho(X,Y)| = 1$ if and only if
$Y = aX + b$ almost surely for some constants $a \neq 0, b$.
\end{proposition}

\begin{proof}
For any $t \in \R$, consider the non-negative quantity
\[
    0 \leq \Var(X + tY)
    = \Var(X) + 2t\,\Cov(X,Y) + t^2 \Var(Y).
\]
This is a quadratic in $t$ with a non-negative minimum, so its
discriminant must be non-positive:
$4\Cov(X,Y)^2 - 4\Var(X)\Var(Y) \leq 0$,
which gives $\Cov(X,Y)^2 \leq \Var(X)\Var(Y)$ and hence
$|\rho(X,Y)| \leq 1$.

Equality $|\rho| = 1$ holds if and only if the discriminant is zero,
meaning the quadratic $\Var(X + tY)$ touches zero at some
$t_0 = -\Cov(X,Y)/\Var(Y)$. Then
$\Var(X + t_0 Y) = 0$, so $X + t_0 Y = c$ almost surely for some
constant $c$, giving $X = -t_0 Y + c$, i.e., $Y = aX + b$ with
$a = -1/t_0 \neq 0$.
\end{proof}

\begin{rlconnection}[title={Variance Reduction in REINFORCE}]
The REINFORCE gradient estimator is
$\hat{g} = \sum_{t=0}^{T-1} \nabla_\theta \log\pi_\theta(A_t \mid S_t)
\cdot (G_t - b(S_t))$,
where $G_t$ is the return and $b(S_t)$ is a baseline. By linearity,
\[
    \E[\nabla\log\pi \cdot b(S)] = \E[\nabla\log\pi] \cdot b(S) = 0
\]
(under regularity conditions), so the baseline does not bias the
estimator. However, by the variance formula,
$\Var(\nabla\log\pi \cdot (G - b))
= \Var(\nabla\log\pi \cdot G) - 2\Cov(\nabla\log\pi \cdot G,
\nabla\log\pi \cdot b) + \Var(\nabla\log\pi \cdot b)$.
Choosing $b$ to minimize this variance (often $b(s) \approx V^\pi(s)$)
is a direct application of the covariance and variance algebra developed
in this section.
\end{rlconnection}

%% ============================================================
%% SECTION 5: JOINT DISTRIBUTIONS
%% ============================================================

\section{Joint Distributions}
\label{sec:joint}

We now formalize the joint behavior of two or more random variables.

\subsection{Joint PMF and PDF}

\begin{definition}[Joint PMF]
\label{def:joint_pmf}
Let $X$ and $Y$ be discrete random variables. The \emph{joint PMF} is
\[
    p_{X,Y}(x, y) = \Prob(X = x, Y = y).
\]
It satisfies $p_{X,Y}(x,y) \geq 0$ for all $(x,y)$ and
$\sum_x \sum_y p_{X,Y}(x,y) = 1$.
\end{definition}

\begin{definition}[Joint PDF]
\label{def:joint_pdf}
Random variables $X$ and $Y$ are \emph{jointly continuous} if there
exists a non-negative function $f_{X,Y} \colon \R^2 \to [0,\infty)$
such that for every region $B \subseteq \R^2$,
\[
    \Prob((X, Y) \in B) = \iint_B f_{X,Y}(x, y) \, dx \, dy.
\]
The function $f_{X,Y}$ is the \emph{joint PDF}, satisfying
$\iint_{\R^2} f_{X,Y}(x,y)\,dx\,dy = 1$.
\end{definition}

\subsection{Marginal Distributions}

\begin{definition}[Marginal Distributions]
\label{def:marginals}
The \emph{marginal PMF} of $X$ from the joint PMF is
$p_X(x) = \sum_y p_{X,Y}(x, y)$.
The \emph{marginal PDF} of $X$ from the joint PDF is
$f_X(x) = \int_{-\infty}^{\infty} f_{X,Y}(x, y) \, dy$.
\end{definition}

\begin{example}[Joint and Marginal PMF]
\label{ex:joint_pmf}
Let $(X, Y)$ have the joint PMF given by the following table:
\[
    \begin{array}{c|cc|c}
    p_{X,Y} & Y = 0 & Y = 1 & p_X \\ \hline
    X = 0 & 0.2 & 0.3 & 0.5 \\
    X = 1 & 0.1 & 0.4 & 0.5 \\ \hline
    p_Y & 0.3 & 0.7 &
    \end{array}
\]
The marginal PMFs are $p_X(0) = 0.5$, $p_X(1) = 0.5$,
$p_Y(0) = 0.3$, $p_Y(1) = 0.7$.
\end{example}

\begin{example}[Joint Uniform Distribution]
\label{ex:joint_uniform}
Let $(X, Y)$ be uniformly distributed on the unit square $[0,1]^2$, so
$f_{X,Y}(x,y) = 1$ for $(x,y) \in [0,1]^2$ and $0$ otherwise. The
marginal of $X$ is
$f_X(x) = \int_0^1 1\,dy = 1$ for $x \in [0,1]$,
i.e., $X \sim \mathrm{Uniform}(0,1)$, and similarly for $Y$.
\end{example}

\subsection{Independence of Random Variables}

\begin{definition}[Independence of Random Variables]
\label{def:rv_independence}
Discrete random variables $X$ and $Y$ are \emph{independent} if
$p_{X,Y}(x, y) = p_X(x) \, p_Y(y)$ for all $x, y$.
Continuous random variables $X$ and $Y$ are \emph{independent} if
$f_{X,Y}(x, y) = f_X(x) \, f_Y(y)$ for all $x, y$.
\end{definition}

\begin{example}[Testing Independence from a Joint PMF]
\label{ex:test_independence}
Continuing Example~\ref{ex:joint_pmf}, check whether $X$ and $Y$ are
independent. We need $p_{X,Y}(x,y) = p_X(x) p_Y(y)$ for all $(x,y)$.
At $(0, 0)$: $p_{X,Y}(0,0) = 0.2$ but
$p_X(0) p_Y(0) = 0.5 \times 0.3 = 0.15 \neq 0.2$.
Therefore $X$ and $Y$ are \emph{not} independent.
\end{example}

\subsection{Conditional Distributions}

\begin{definition}[Conditional PMF]
\label{def:cond_pmf}
If $X$ and $Y$ are discrete with $p_Y(y) > 0$, the \emph{conditional
PMF of $X$ given $Y = y$} is
\[
    p_{X \mid Y}(x \mid y)
    = \frac{p_{X,Y}(x, y)}{p_Y(y)}.
\]
\end{definition}

\begin{definition}[Conditional PDF]
\label{def:cond_pdf}
If $X$ and $Y$ are jointly continuous with $f_Y(y) > 0$, the
\emph{conditional PDF of $X$ given $Y = y$} is
\[
    f_{X \mid Y}(x \mid y)
    = \frac{f_{X,Y}(x, y)}{f_Y(y)}.
\]
\end{definition}

\begin{example}[Conditional PMF Computation]
\label{ex:cond_pmf_compute}
From Example~\ref{ex:joint_pmf}, the conditional PMF of $X$ given
$Y = 1$ is
\[
    p_{X \mid Y}(0 \mid 1) = \frac{0.3}{0.7} = \frac{3}{7}, \qquad
    p_{X \mid Y}(1 \mid 1) = \frac{0.4}{0.7} = \frac{4}{7}.
\]
Given $Y = 1$, the conditional distribution of $X$ is
$\mathrm{Bernoulli}(4/7)$.
\end{example}

%% ============================================================
%% SECTION 6: CONDITIONAL EXPECTATION AND THE TOWER PROPERTY
%% ============================================================

\section{Conditional Expectation and the Tower Property}
\label{sec:cond_exp}

\begin{definition}[Conditional Expectation Given an Event]
\label{def:cond_exp_event}
Let $X$ be a random variable and $B$ an event with $\Prob(B) > 0$. The
\emph{conditional expectation of $X$ given $B$} is
\[
    \E[X \mid B] = \sum_x x \, \Prob(X = x \mid B)
    \quad \text{(discrete)}
\]
or $\E[X \mid B] = \int x \, f_{X \mid B}(x)\,dx$ (continuous).
\end{definition}

\begin{definition}[Conditional Expectation Given a Random Variable]
\label{def:cond_exp_rv}
Let $X$ and $Y$ be random variables. The \emph{conditional expectation
of $X$ given $Y = y$} is
\[
    \E[X \mid Y = y] = \begin{cases}
    \sum_x x \, p_{X \mid Y}(x \mid y) & \text{(discrete)}, \\
    \int_{-\infty}^{\infty} x \, f_{X \mid Y}(x \mid y)\,dx
    & \text{(continuous)}.
    \end{cases}
\]
The \emph{conditional expectation of $X$ given $Y$}, denoted
$\E[X \mid Y]$, is the random variable obtained by evaluating
$\E[X \mid Y = y]$ at $Y$: it equals $h(Y)$ where $h(y) = \E[X \mid Y = y]$.
\end{definition}

\begin{keyresult}[title={Law of Total Expectation (Tower Property)}]
The law of total expectation, also called the tower property or iterated
expectation, states that $\E[\E[X \mid Y]] = \E[X]$. This identity is
the algebraic backbone of the Bellman equation: by conditioning on the
next state $S_{t+1}$, the value function at time $t$ decomposes into an
immediate reward and a conditional future value.
\end{keyresult}

\begin{theorem}[Law of Total Expectation]
\label{thm:tower}
Let $X$ and $Y$ be random variables with $\E[|X|] < \infty$. Then
\[
    \E\bigl[\E[X \mid Y]\bigr] = \E[X].
\]
\end{theorem}

\begin{proof}
We prove the discrete case; the continuous case is analogous with
integrals replacing sums.

Let $h(y) = \E[X \mid Y = y] = \sum_x x \, p_{X \mid Y}(x \mid y)$.
Then $\E[X \mid Y] = h(Y)$ is a function of $Y$, so
\begin{align*}
    \E\bigl[\E[X \mid Y]\bigr]
    &= \E[h(Y)]
    = \sum_y h(y) \, p_Y(y) \\
    &= \sum_y \left(\sum_x x \, p_{X \mid Y}(x \mid y)\right) p_Y(y) \\
    &= \sum_y \sum_x x \, \frac{p_{X,Y}(x,y)}{p_Y(y)} \, p_Y(y) \\
    &= \sum_y \sum_x x \, p_{X,Y}(x,y)
    = \sum_x x \sum_y p_{X,Y}(x,y)
    = \sum_x x \, p_X(x) = \E[X].
\end{align*}
The key step is that $p_{X\mid Y}(x \mid y) \cdot p_Y(y) = p_{X,Y}(x,y)$
cancels the denominator, recovering the joint PMF.
\end{proof}

\begin{example}[Tower Property in Action]
\label{ex:tower_concrete}
Roll a fair die ($Y$), then toss a fair coin $Y$ times. Let $X$ be the
number of heads. Given $Y = y$, we have
$X \mid (Y = y) \sim \mathrm{Bin}(y, 1/2)$, so
$\E[X \mid Y = y] = y/2$, and thus $\E[X \mid Y] = Y/2$.
By the tower property,
$\E[X] = \E[\E[X \mid Y]] = \E[Y/2] = \E[Y]/2 = 3.5/2 = 1.75$.
\end{example}

\begin{rlconnection}[title={Tower Property and the Bellman Equation}]
The Bellman equation for the state-value function is derived by
conditioning on the first transition:
\begin{align*}
    V^\pi(s) &= \E\!\left[\sum_{t=0}^\infty \gamma^t R_t
    \;\middle|\; S_0 = s\right]
    = \E\!\left[R_0 + \gamma \sum_{t=1}^\infty \gamma^{t-1} R_t
    \;\middle|\; S_0 = s\right] \\
    &= \E[R_0 \mid S_0 = s]
    + \gamma \, \E\!\left[\E\!\left[\sum_{t=1}^\infty
    \gamma^{t-1} R_t \;\middle|\; S_1\right]
    \;\middle|\; S_0 = s\right] \\
    &= \E[R_0 \mid S_0 = s] + \gamma \, \E[V^\pi(S_1) \mid S_0 = s].
\end{align*}
The second line applies the tower property: the outer expectation is over
$S_1$ (given $S_0 = s$), and the inner expectation is over the future
trajectory given $S_1$. Without the tower property, the recursive
decomposition of value functions would not be justified.
\end{rlconnection}

\begin{example}[Conditional Expectation with Continuous Variables]
\label{ex:cond_exp_continuous}
Let $X \sim \mathrm{Exp}(1)$ and $Y \mid (X = x) \sim
\mathrm{Uniform}(0, x)$. Then
$\E[Y \mid X = x] = x/2$, so $\E[Y \mid X] = X/2$.
By the tower property,
$\E[Y] = \E[X/2] = \E[X]/2 = 1/2$.
\end{example}

\subsection{Law of Total Variance}

\begin{theorem}[Law of Total Variance]
\label{thm:total_variance}
Let $X$ and $Y$ be random variables with $\E[X^2] < \infty$. Then
\[
    \Var(X) = \E[\Var(X \mid Y)] + \Var(\E[X \mid Y]),
\]
where $\Var(X \mid Y)$ is the random variable
$\E[X^2 \mid Y] - (\E[X \mid Y])^2$.
\end{theorem}

\begin{proof}
By the computational formula for variance,
$\Var(X) = \E[X^2] - (\E[X])^2$.
We compute each term using the tower property.

\emph{First term.} By the tower property,
$\E[X^2] = \E[\E[X^2 \mid Y]]$.

\emph{Second term.} By the tower property,
$\E[X] = \E[\E[X \mid Y]]$, so $(\E[X])^2 = (\E[\E[X \mid Y]])^2$.

Now consider the right-hand side. The conditional variance is
$\Var(X \mid Y) = \E[X^2 \mid Y] - (\E[X \mid Y])^2$. Taking
expectations:
\[
    \E[\Var(X \mid Y)] = \E[\E[X^2 \mid Y]] - \E[(\E[X \mid Y])^2]
    = \E[X^2] - \E[(\E[X \mid Y])^2].
\]
The variance of the conditional expectation is
\[
    \Var(\E[X \mid Y]) = \E[(\E[X \mid Y])^2] - (\E[\E[X \mid Y]])^2
    = \E[(\E[X \mid Y])^2] - (\E[X])^2.
\]
Adding:
\begin{align*}
    \E[\Var(X \mid Y)] + \Var(\E[X \mid Y])
    &= \bigl(\E[X^2] - \E[(\E[X \mid Y])^2]\bigr)
    + \bigl(\E[(\E[X \mid Y])^2] - (\E[X])^2\bigr) \\
    &= \E[X^2] - (\E[X])^2 = \Var(X). \qedhere
\end{align*}
\end{proof}

\begin{example}[Total Variance with Die and Coin]
\label{ex:total_variance}
Continuing Example~\ref{ex:tower_concrete}: $Y$ is a die roll,
$X \mid (Y = y) \sim \mathrm{Bin}(y, 1/2)$.
\begin{itemize}[nosep]
    \item $\E[X \mid Y] = Y/2$, so $\Var(\E[X \mid Y])
    = \Var(Y/2) = \Var(Y)/4 = (35/12)/4 = 35/48$.
    \item $\Var(X \mid Y) = Y \cdot (1/2)(1/2) = Y/4$, so
    $\E[\Var(X \mid Y)] = \E[Y]/4 = 3.5/4 = 7/8$.
    \item $\Var(X) = 7/8 + 35/48 = 42/48 + 35/48 = 77/48 \approx 1.604$.
\end{itemize}
The within-group variance ($7/8$) and the between-group variance
($35/48$) together account for the total.
\end{example}

%% ============================================================
%% SECTION 7: EXERCISES
%% ============================================================

\section{Exercises}
\label{sec:exercises}

\begin{exercise}[Indicator Expectation and Probability]
\textnormal{(\(\ast\))} \\
Let $A_1, \ldots, A_n$ be events. Using the indicator random variables
$\indicator_{A_i}$ and linearity of expectation, prove that
\[
    \Prob\!\left(\bigcup_{i=1}^n A_i\right)
    \leq \sum_{i=1}^n \Prob(A_i)
\]
(the union bound). Compare this proof with the inclusion-exclusion
approach from Lesson~9.
\end{exercise}

\begin{exercise}[Variance of a Linear Combination]
\textnormal{(\(\ast\))} \\
Let $X_1, X_2, X_3$ be random variables with $\Var(X_i) = \sigma^2$
for each $i$ and $\Cov(X_i, X_j) = \rho \sigma^2$ for $i \neq j$,
where $\rho \in [-1/2, 1]$.
\begin{enumerate}[nosep,label=(\alph*)]
    \item Compute $\Var(X_1 + X_2 + X_3)$ in terms of $\sigma^2$ and
    $\rho$.
    \item For what value of $\rho$ is the variance of the sum minimized?
    What is the minimum?
\end{enumerate}
\end{exercise}

\begin{exercise}[Conditional Expectation Practice]
\textnormal{(\(\ast\ast\))} \\
A box contains $N$ balls numbered $1, \ldots, N$, where $N$ itself is
random with $\Prob(N = n) = 1/4$ for $n = 1, 2, 3, 4$. A ball is drawn
uniformly at random from the box; let $X$ be its number.
\begin{enumerate}[nosep,label=(\alph*)]
    \item Compute $\E[X \mid N = n]$ for each $n$.
    \item Use the tower property to compute $\E[X]$.
    \item Use the law of total variance to compute $\Var(X)$.
\end{enumerate}
\end{exercise}

\begin{exercise}[Uncorrelated but Dependent -- Discrete Case]
\textnormal{(\(\ast\ast\))} \\
Let $Y$ be uniformly distributed on $\{-1, 0, 1\}$ and define
$X = Y^2$.
\begin{enumerate}[nosep,label=(\alph*)]
    \item Compute $\E[X]$, $\E[Y]$, $\E[XY]$, and $\Cov(X, Y)$.
    \item Verify that $X$ and $Y$ are uncorrelated.
    \item Show that $X$ and $Y$ are \emph{not} independent by finding
    values $x, y$ such that $p_{X,Y}(x, y) \neq p_X(x) p_Y(y)$.
\end{enumerate}
\end{exercise}

\begin{exercise}[Tower Property and Bellman Equation]
\textnormal{(\(\ast\ast\ast\))} \\
Consider an MDP with states $\{s_1, s_2\}$, actions $\{a\}$ (single
action), discount factor $\gamma = 0.9$, reward
$R(s_1) = 1$, $R(s_2) = 0$, and transition probabilities
$\Prob(S_{t+1} = s_1 \mid S_t = s_1) = 0.8$ and
$\Prob(S_{t+1} = s_1 \mid S_t = s_2) = 0.4$.
\begin{enumerate}[nosep,label=(\alph*)]
    \item Write out the Bellman equations for $V(s_1)$ and $V(s_2)$
    using the tower property decomposition.
    \item Solve the system of linear equations for $V(s_1)$ and
    $V(s_2)$.
    \item Verify your answer by computing
    $V(s_1) = \E[R_0 \mid S_0 = s_1] + \gamma \, \E[V(S_1) \mid
    S_0 = s_1]$ explicitly.
\end{enumerate}
\end{exercise}

\begin{exercise}[Law of Total Variance for Mixture Distributions]
\textnormal{(\(\ast\ast\ast\))} \\
A reward distribution is a mixture: with probability $\alpha$, the
reward is $R \sim \mathcal{N}(\mu_1, \sigma_1^2)$, and with probability
$1 - \alpha$, the reward is $R \sim \mathcal{N}(\mu_2, \sigma_2^2)$.
Let $Z \in \{1, 2\}$ indicate the component.
\begin{enumerate}[nosep,label=(\alph*)]
    \item Use the tower property to compute $\E[R]$.
    \item Use the law of total variance to compute $\Var(R)$ and show
    that it decomposes as
    \[
        \Var(R) = \underbrace{\alpha\sigma_1^2 +
        (1-\alpha)\sigma_2^2}_{\text{within-component}}
        + \underbrace{\alpha(1-\alpha)(\mu_1 - \mu_2)^2}_{%
        \text{between-component}}.
    \]
    \item Explain why mixture distributions always have higher variance
    than any single component when $\mu_1 \neq \mu_2$.
\end{enumerate}
\textit{Hint:} For part (b), compute $\Var(\E[R \mid Z])$ by noting
that $\E[R \mid Z]$ takes two values with known probabilities.
\end{exercise}

%% ============================================================
%% SECTION 8: SUMMARY
%% ============================================================

\section{Summary}
\label{sec:summary}

\begin{itemize}
    \item The \emph{Law of the Unconscious Statistician} (LOTUS) states
    $\E[g(X)] = \sum_x g(x) p_X(x)$ (discrete) or
    $\E[g(X)] = \int g(x) f_X(x)\,dx$ (continuous), allowing
    computation of expectations of transformed variables without
    finding the new distribution.
    \item \emph{Linearity of expectation}
    ($\E[aX + bY] = a\E[X] + b\E[Y]$) holds for \emph{all} random
    variables with finite expectations, with \emph{no independence
    assumption required}.
    \item Every probability can be expressed as the expectation of an
    indicator: $\Prob(A) = \E[\indicator_A]$.
    \item The \emph{variance} satisfies $\Var(aX + b) = a^2\Var(X)$
    and the computational formula
    $\Var(X) = \E[X^2] - (\E[X])^2$.
    \item \emph{Covariance} $\Cov(X,Y) = \E[XY] - \E[X]\E[Y]$ is
    a symmetric, bilinear measure of linear association.
    \item The \emph{variance of a sum} decomposes as
    $\Var(\sum X_i) = \sum \Var(X_i) + 2\sum_{i<j} \Cov(X_i,X_j)$;
    for independent variables, the cross-covariance terms vanish.
    \item Independence implies zero covariance, but \emph{the converse
    is false}: uncorrelated variables can be strongly dependent
    (e.g., $X$ and $X^2$ when $X$ is symmetric around zero).
    \item The \emph{conditional expectation} $\E[X \mid Y]$ is a
    random variable (a function of $Y$) that satisfies the
    \emph{tower property}: $\E[\E[X \mid Y]] = \E[X]$.
    \item The \emph{law of total variance}
    $\Var(X) = \E[\Var(X \mid Y)] + \Var(\E[X \mid Y])$ decomposes
    total variability into within-group and between-group components.
\end{itemize}

%% ============================================================
%% SECTION 9: FURTHER READING
%% ============================================================

\section{Further Reading}
\label{sec:further_reading}

\begin{itemize}
    \item \textbf{Lesson 12} develops the law of large numbers and the
    central limit theorem, which explain how sample means concentrate
    and why the normal distribution appears as a universal limit.
    \item \textbf{Phase 01, Lesson 5} (probability foundations) revisits
    expectation and conditional expectation in the measure-theoretic
    framework, defining $\E[X \mid \mathcal{G}]$ as a
    $\mathcal{G}$-measurable random variable satisfying a projection
    property. The tower property becomes a consequence of the
    $\sigma$-algebra nesting $\mathcal{G}_1 \subseteq \mathcal{G}_2$.
    \item \textbf{Phase 02} (Temporal Difference Learning, Policy
    Gradient Theorem) relies heavily on the linearity of expectation
    and the tower property for deriving TD update rules and the
    policy gradient theorem.
    \item For additional worked examples on conditional expectation and
    the tower property, see \citet[Chapter 9]{blitzstein2019}, which
    includes the ``conditional expectation as a projection''
    perspective.
    \item The law of total variance is sometimes called the
    \emph{Eve's law} (``E of the V plus V of the E''); a thorough
    discussion with applications to hierarchical models appears in
    \citet[Section 4.4]{bertsekas2008}.
    \item The Cauchy-Schwarz inequality for random variables
    ($|\E[XY]|^2 \leq \E[X^2]\E[Y^2]$), which underlies the
    correlation bound $|\rho| \leq 1$, has far-reaching consequences
    in functional analysis; see Phase~01 (Lesson~7, inner product
    spaces).
\end{itemize}

\bibliography{probability-refs}

\end{document}
